\chapter{Methodology}

This chapter sets out the method used to evaluate the prototype. The study was conceived as a small-scale exploratory evaluation of a deliberately scoped educational system; it was not designed as a controlled comparison or as a test of whether VR should replace conventional teaching.

\section{Research Framing}

The evaluation considered whether the prototype was educationally useful within its intended scope, whether participants could engage with it clearly enough for that value to matter, and whether the combined evidence justified further development for selected computer architecture concepts.

As noted in Chapter~1, the prototype grew out of repeated teaching-assistant experience in computer organisation. The study therefore examined whether the guided tracing that proved useful in that setting could be translated into a useful VR-based learning environment, a question that also sits in line with earlier architecture-focused immersive work \cite{dascalu2017experiential,tan2023vr}.

\subsection{Research Question}

The dissertation was guided by the following research question:

\begin{quote}
\textit{To what extent can an immersive virtual reality learning environment support students' understanding of selected single-cycle MIPS datapath concepts as a supplementary educational tool?}
\end{quote}

In this dissertation, the phrase ``to what extent'' is to be interpreted in a bounded exploratory sense rather than as a promise of a full classroom control-group comparison. More specifically, the question is answered through three linked forms of evidence: participant perceptions of the environment, descriptive concept evidence drawn from the baseline subset and post-session responses, and observed progression or support needs across the session structure.

The purpose of the study was therefore to determine whether this particular prototype, used in this particular educational context, offered enough conceptual clarity, usability, and perceived value to justify the design choices behind it and support later work on broader evaluation.

\subsection{Evaluation Aspects}

The research question was examined through four linked evaluation aspects reflecting the prototype's intended educational role:

\begin{enumerate}
\item \textbf{Usability and clarity.} Whether participants could understand, navigate, and use the system clearly enough for its educational purpose to be evaluated.
\item \textbf{Conceptual support.} Whether the collected evidence suggested that the environment helped participants follow instruction flow, datapath phases, component roles, and selected architectural relationships.
\item \textbf{Progression from guided to reduced support.} Whether participants could move from the guided structure of Learning mode into less-supported interaction in Practice or Test mode, and what support remained necessary during that transition.
\item \textbf{Supplementary educational value.} How participants perceived the prototype in relation to lectures, diagrams, tutorials, simulators, and other conventional teaching materials.
\end{enumerate}

\section{Alignment with the Learning Outcomes}

The evaluation was aligned with the learning outcomes defined for the prototype, but those outcomes belong primarily to the design of the system and are discussed in full in Section~4.1. In methodological terms, they matter here because they establish what the study should look for and what kinds of claims the evidence can reasonably support.

The study therefore focused on a bounded cluster of educational tasks: decoding supported instructions, explaining the roles of major datapath components, tracing instruction and value flow across the main lesson phases, making the required operand and control-facing choices, predicting resulting architectural state changes, reasoning across different supported instruction classes, and completing the walkthrough with progressively less support. These categories were reflected across the baseline subset, the post-session knowledge check, the observation notes, and the broader interpretation of participant feedback.

\section{Study Design and Participants}

Structured evidence came from the background form, the baseline subset, the Likert-style post-session items addressing perceptions of the system, and the knowledge check on selected concepts. Open-ended evidence came from written participant feedback and researcher observation notes recorded during and after the session. These sources were considered together to provide a fuller account of how the prototype functioned as a learning tool.

The evaluation remained exploratory rather than classroom-trial in scale. There was no formal non-VR control group and the study did not attempt to measure long-term retention. The design instead focused on how participants engaged with the prototype and whether the collected evidence supported its intended role as a supplementary learning environment, consistent with broader literature on conceptual VR support in computer science education \cite{janani2026virtual,dengel2020important}.

Twelve participants took part in the completed study through voluntary convenience sampling. All participants were consenting adults aged 18 or over. The approved ethics application had initially set a target sample size of twenty participants, while the completed recruitment yielded twelve. Recruitment took place by email invitation, no payment was offered, and the study involved no dependent relationship between the research team and participants. Background information on prior architecture exposure, prior VR use, confidence, and topic familiarity was collected before the session so that the findings could be interpreted in light of participants' different starting points; these characteristics are reported in Chapter~6.

\section{Procedure}

Each participant session followed the same broad structure: consent and background capture, a short baseline probe, a tailored refresher, the VR session itself, and a post-session evaluation.

The session began with the Participant Information Leaflet and written informed consent, both of which are reproduced in Appendix~\ref{app:pil} and Appendix~\ref{app:consent} respectively. Participants were given the opportunity to read the study information, ask questions, and decide whether they wished to continue before any data collection began. They then completed the background section of the questionnaire reproduced in Appendix~\ref{app:questionnaire}, which captured prior VR use, prior exposure to computer architecture or computer organisation, and self-reported familiarity and confidence.

After the background section, participants were offered a short baseline subset drawn from the knowledge check reproduced in Appendix~\ref{app:questionnaire}. This consisted of five pre-session questions used as a modest checkpoint on foundational ideas before the VR experience began, specifically Questions 11--14 and 25 from the post-session questionnaire in Appendix~\ref{app:questionnaire}. Because this baseline acted as a quick pre-session scan rather than a mandatory scored gate, unanswered items were retained as missing rather than treated as incorrect. Participants then received a short verbal refresher tailored to their prior experience. The refresher was intentionally high-level: it was not meant to teach the whole content in advance, but to align terminology, reduce avoidable confusion, and give participants enough orientation to engage meaningfully with the prototype. Typical coverage included the basic form of an assembly instruction, what the datapath represents, what the main instruction stages do, how registers and memory differ, where control choices matter, and what kinds of interaction participants would encounter in the VR environment.

Participants then completed the VR session using selected instructions and one or more of the available modes. A recorded walkthrough of the prototype, representative of what participants were seeing during the live study, is available at \url{https://youtu.be/Y8tFYnWeLtI}. The exact path depended on participant background, session time, and whether the participant appeared comfortable enough to move into reduced-guidance modes. A typical path included one or two Learning mode instructions, one Practice mode instruction, and, where appropriate, a further Test mode attempt. During the session, the researcher intervened only when needed, and participants were encouraged to make full use of the resources available inside the environment. The amount and kind of intervention required were themselves treated as meaningful observational data.

After removing the headset, participants completed the remaining sections of the questionnaire reproduced in Appendix~\ref{app:questionnaire}. These included Likert-style ratings on usability, engagement, and perceived learning, followed by the broader knowledge-check section and open-ended feedback. Alongside the formal questionnaire data, the researcher maintained brief session notes on participant support, attempted instructions and modes, points of hesitation, and relevant comparisons with traditional materials.

\section{Instruments and Data Capture}

The primary structured instrument was the Microsoft Forms questionnaire used in the live study. It captured participant background, post-session usability, engagement and perceived-learning responses, the targeted knowledge check and follow-up confidence item, and open-ended feedback. A reconstructed print version of the deployed questionnaire is reproduced in Appendix~\ref{app:questionnaire}, while the Participant Information Leaflet, consent form, and representative recruitment materials are reproduced in Appendix~\ref{app:pil}, Appendix~\ref{app:consent}, and Appendix~\ref{app:recruitment}.

The knowledge-check portion of the live form focused mainly on topics relevant to the implemented prototype: program-counter behaviour, fetch, the instruction register, stage order, the execute stage, branch-related reasoning, and the behaviour of supported MIPS instructions such as \texttt{add}, \texttt{addi}, \texttt{lw}, \texttt{sw}, \texttt{beq}, and \texttt{j}. A small number of surrounding architecture concepts were also retained so that responses could be interpreted against the broader background knowledge with which participants entered the session, while the main evaluative focus remained on the selected datapath concepts supported by the prototype.

Some participant-facing materials, most explicitly the Participant Information Leaflet in Appendix~\ref{app:pil}, described the study using broader computer architecture examples such as pipelining, memory hierarchy, and stack behaviour. These materials are reproduced unchanged for transparency because they formed part of the approved ethics documentation before the final prototype scope was fully settled. The evaluated VR session itself remained limited to the selected single-cycle MIPS datapath concepts defined in this dissertation.

Researcher notes supplemented the questionnaire by recording attempted instructions and modes, points of hesitation, support requirements, and relevant comparisons participants made with conventional learning materials.

\section{Analysis Plan}

The analysis was primarily descriptive and combined evidence on usability, conceptual support, participant understanding, and supplementary educational value. The small sample and exploratory design did not justify strong inferential testing, so the evidence was interpreted through descriptive statistics, item-level patterns, open-ended feedback, and researcher observations. IBM SPSS Statistics was used for quantitative analysis, while NVivo supported the organisation of qualitative material.

Quantitative analysis summarised participant background, item-level Likert responses and descriptive group measures, knowledge-check performance, post-knowledge-check confidence, and the limited baseline/post-session comparison. Frequencies, percentages, means, medians, ranges, and internal-consistency estimates were reported where appropriate to the measure being described. Baseline and post-session responses were compared descriptively only on the repeated items for which matched data were available.

Open-ended responses and researcher observation notes were read together through a light thematic analysis. The analysis focused on recurring material concerning what helped understanding, what remained confusing, interface or navigation friction, the guidance structure, and comparisons with conventional learning materials.

Because the study did not include a formal non-VR control group, comparison remained indirect and drew on participant perceptions, the limited repeated baseline/post-session items, and observed progression from guided to less-guided performance. The baseline subset was followed by a short verbal refresher before the VR session itself, so any baseline/post-session differences were interpreted as performance following the overall learning session rather than as evidence attributable to the VR environment alone.

\section{Ethical Considerations and Study Boundaries}

The study received research ethics approval through the School of Computer Science and Statistics Research Ethics Management System (REAMS). Participation was voluntary and based on informed consent, with participants receiving study information before taking part and retaining the right to withdraw without penalty. They were also informed that withdrawal of research material might no longer be possible once data had been anonymised and incorporated into the analysis.

The principal participant risk was temporary discomfort associated with VR use, including eye strain, dizziness, or motion sickness. Participants could pause or stop the session immediately if they felt unwell, and task completion was not prioritised over participant comfort.

Administrative contact information was kept separate from research responses, and research materials were anonymised as early as practicable. Access to identifiable data was restricted to what was necessary for the study, with storage limited to password-protected and university-approved systems in line with the approved data-handling plan. Identifiable or coded administrative data were retained only for a limited period before deletion.

The study used a Meta VR headset supplied through the School, and participants were not required to sign into personal Meta accounts. The research team did not intentionally collect Meta account data, although participant materials noted that limited device-level or platform-level processing outside the direct control of the research team could still occur.

\section{Use of Generative Artificial Intelligence}

In accordance with Trinity College Dublin's \textit{College Statement on Artificial Intelligence and Generative AI in Teaching, Learning, Assessment \& Research} (2024), the use of generative artificial intelligence in the development of this project and preparation of the dissertation is acknowledged below.

ChatGPT (GPT-5.6 Sol, OpenAI, \url{https://chatgpt.com}) was used in a supporting capacity to clarify selected computer-architecture concepts during development, including walkthroughs of how supported instructions operate through a single-cycle datapath. It was also used to help standardise BibTeX entries where citation information obtained from Google Scholar and individual publisher websites was inconsistent, and to provide critique and proofreading of author-written dissertation material.

Codex (GPT-5.4, OpenAI, \url{https://openai.com/codex/}) was used in a limited support capacity. Its primary use was to help maintain and organise a project journal so that implementation milestones and decisions could later be referred back to as historical summaries. It also provided limited assistance with repository organisation, the Python data-processing pipeline used to clean, restructure, and prepare collected study data for analysis, and the preparation of tables. Codex was not used to create, fabricate, invent, simulate, or alter participant data, results, or findings.

\section{Summary}

This chapter established the methodological framework for the exploratory evaluation, including the study design and participants, procedure, data-collection instruments, analysis strategy, and ethical considerations. The analysis remained primarily descriptive, while the tailored refresher between the baseline probe and VR session means that any baseline/post-session differences reflect the overall learning session rather than an effect attributable to VR alone. Chapter~4 turns from this methodological framework to the design decisions through which the targeted datapath concepts were translated into a guided, learner-facing VR environment.